% The bring-up order, and the three checks that stop the chapter rather than
% letting it continue with a fault.
\begin{tikzpicture}[font=\sffamily\scriptsize, x=1cm, y=1cm]

\tikzset{
  stage/.style={draw=linecol, fill=kercol, rounded corners=2pt, align=center,
                text width=34mm, minimum height=10mm, inner sep=3pt},
  gate/.style={draw=red!70!black, fill=red!8, rounded corners=2pt, align=left,
               text width=48mm, inner sep=3pt},
  needs/.style={font=\sffamily\tiny, text=black!55, align=left},
}

\foreach \i/\t/\n in {%
  0/{1. select the kernel clock\\{\tiny and read it back from the registers}}/{nothing},
  1/{2. compute both phases\\{\tiny nominal and data, with their sample points}}/{nothing},
  2/{3. lay out the message memory\\{\tiny filters, receive and transmit buffers}}/{nothing},
  3/{4. internal loopback\\{\tiny all three frame formats, payload compared}}/{nothing},
  4/{5. delay compensation\\{\tiny from the transceiver's loop delay}}/{the datasheet},
  5/{6. external loopback\\{\tiny proves the pins and the alternate functions}}/{one wire},
  6/{7. the bus\\{\tiny transceiver, termination, and chapter 10's master}}/{a purchase}}{
  \node[stage] (s\i) at (0,-\i*1.5) {\t};
  \node[needs, anchor=west] at (2.1,-\i*1.5) {needs: \n};
}
\foreach \i in {0,1,2,3,4,5}{ \draw[flow] (s\i) -- (s\the\numexpr\i+1\relax); }

% ---------------- the three checks that stop the node
\node[gate] (g0) at (-5.4,0)    {\textbf{stop} if no kernel clock source is selected. Every number below this depends on it};
\node[gate] (g1) at (-5.4,-1.5) {\textbf{stop} if either phase does not divide exactly. A rate that is one part in a thousand out works only against a node that is wrong the same way};
\node[gate] (g2) at (-5.4,-3.0) {\textbf{stop} at build time if the layout does not fit the part's message memory};

\draw[faultedge] (s0.west) -- (g0.east);
\draw[faultedge] (s1.west) -- (g1.east);
\draw[faultedge] (s2.west) -- (g2.east);

\node[umlnote, text width=52mm, anchor=north west] at (-8.0,-5.2)
  {The first four stages need no hardware beyond the board, which is why they
   come first: the bit timing arithmetic, the memory layout and all three frame
   formats can be proven in internal loopback while the transceiver is still in
   the post.};
\node[umlnote, text width=52mm, anchor=north west] at (-8.0,-7.6)
  {The last stage is the one that produces the most confusing symptom in the
   whole peripheral: a node alone on a terminated bus retransmits a frame nobody
   will acknowledge, the error counters rise, and from the firmware's side it
   looks exactly like a transmitter that is not working. Chapter 12 makes that
   behaviour the subject rather than the surprise.};
\end{tikzpicture}
